% Paper 3, §11 — CoT faithfulness: does the trace cause the answer?
% Anchors: kb/notes/reasoning/chain-of-thought.md §4 faithfulness
%          kb/notes/reasoning/test-time-compute.md §6 interpretability-questions
%          kb/excerpts/{arcuschin2025-cot-wild,shen2025-faithcot-bench,
%                       mehta2026-faithfulness-scaling}.md (excerpts skeleton-only
%                       as of 2026-05; \deepencite where load-bearing).
% Forward refs: §8 R1 (sec:r1-family), §12 PRM survival (sec:prm-survival),
%               §14 contradictions catalogue (sec:contradictions).
% Cross-paper: Paper 4 §interpretability uses CoT as a probe surface;
%              Paper 5 §alignment treats CoT as a monitorable channel.

\section{CoT faithfulness: does the trace cause the answer?}
\label{sec:cot-faithfulness}

The decoding factorisation $z \sim p_\theta(\cdot \mid x)$,
$y \sim p_\theta(\cdot \mid x, z)$ that
\cref{eq:cot-decode} introduces is a generative model: the trace is
sampled, then the answer is sampled conditional on the trace. As a
\emph{statistical} dependence claim this is uncontroversial. The
\emph{causal} claim — that perturbing $z$ propagates to a different $y$
— is not the same statement, and the 2025--2026 wave of empirical work
finds that the gap between the two can be substantial. \textbf{CoT
faithfulness} is the operational name for the question: does the trace
$z$ in fact carry the information that determines $y$, or has $y$
already been chosen by some other mechanism with $z$ generated as a
post-hoc rationalisation that happens to be human-readable? This
section makes the question precise, summarises the three load-bearing
2025--2026 measurements (Arcuschin et al.\ ``CoT in the wild'', Shen et
al.\ FaithCoT-Bench, Mehta et al.\ on whether scaling improves
faithfulness), states the open contradiction with R1-style
reasoning-training framings (\cref{sec:r1-family}), and draws the
cross-paper threads: Paper~4's interpretability programme uses CoT as a
probe surface and Paper~5's alignment programme treats CoT as a
monitorable channel — both are compromised to the extent that traces
are unfaithful.

\subsection{Operational definition: from conditional to causal}
\label{sec:faith-def}

Following the protocol shared across
\citet{arcuschin2025-cot-wild,shen2025-faithcot-bench,
mehta2026-faithfulness-scaling} and the KB note's framing,%
\footnote{KB note:
\texttt{kb/notes/reasoning/chain-of-thought.md\#4}.} the trace $z$ is
\emph{faithful} to the answer $y$ on input $x$ when an intervention on
$z$ that should change the answer in fact does, and an intervention
that should not does not. Formally, fix $x$ and let
$\mathcal{I}: z \mapsto z'$ be an intervention map. The trace is
faithful with respect to $\mathcal{I}$ when
\begin{equation}
  p_\theta\!\bigl(y \mid x, z\bigr) \;\neq\; p_\theta\!\bigl(y \mid x, z' = \mathcal{I}(z)\bigr)
  \label{eq:faith-causal}
\end{equation}
under interventions $\mathcal{I}$ that change the load-bearing content
of the trace (e.g., flipping an arithmetic step, contradicting an
intermediate conclusion, deleting the chain), and is faithful in the
opposite direction when \cref{eq:faith-causal} fails to hold under
interventions that should be answer-irrelevant (e.g., paraphrasing the
trace into different words but the same content). The discrepancy
between $p_\theta(y \mid x, z)$ and $p_\theta(y \mid x, z')$ — measured
by an answer-distribution distance such as the answer-marginal total
variation or the binary indicator $\mathbb{1}[\hat y(z) \neq \hat y(z')]$
— is the empirical surface of the faithfulness question.

Three intervention families dominate the 2025--2026 literature
\citep[\S 3]{shen2025-faithcot-bench}:%
\footnote{Excerpt status:
\texttt{kb/excerpts/shen2025-faithcot-bench.md} is abstract-only as of
2026-05; \deepencite{shen2025-faithcot-bench §3, citation-pending}}
\begin{enumerate}
\item \textbf{Paraphrase $z'$.} Same content, different surface form. A
faithful trace channel should leave the answer distribution unchanged;
the failure mode is the model committing to a different $y$ when an
incidental phrase in $z$ is swapped.
\item \textbf{Contradict $z'$.} Replace one or more steps with an
explicit contradiction (e.g., flip an inequality, negate a numeric
claim). A faithful trace channel should propagate the contradiction
into a changed $y$; the failure mode is the answer surviving an
internally-inconsistent trace.
\item \textbf{Truncate $z'$.} Drop the tail of $z$ before the trace
reaches its conclusion. A faithful trace channel should yield a
substantially worse $y$; the failure mode is the answer being recoverable
from the prompt $x$ alone, with $z$'s contribution to $p_\theta(y \mid
x, z)$ approximating $p_\theta(y \mid x)$.
\end{enumerate}
The third intervention is the cleanest operational test. If
$\mathcal{D}\!\bigl(p_\theta(y \mid x, z),\, p_\theta(y \mid x)\bigr)$
— some divergence between the trace-conditioned and trace-free answer
distributions — is small, the trace is \emph{not contributing} to the
answer, regardless of how reasonable it reads.

\begin{intuition}
Faithfulness is a Granger-causality-like test on the trace channel:
does conditioning on $z$ \emph{change} our prediction of $y$ relative
to conditioning on $x$ alone, and do interventions on $z$ propagate
to $y$ in the directions content-causality would predict? Returning
to math: \cref{eq:faith-causal} compares two answer distributions
indexed by counterfactual traces $z$ and $z' = \mathcal{I}(z)$, and
the operational metric is a divergence between them. The Granger
analogy is loose — Granger causality is a statement about temporal
predictive precedence in linear time series, not about generative
counterfactuals — but the canonical machinery is the same: hold the
prompt fixed, vary the candidate cause, and measure the response in
the answer distribution.
\end{intuition}

\subsection{Arcuschin et al.: CoT in the wild is not always faithful}
\label{sec:faith-arcuschin}

\citet{arcuschin2025-cot-wild} run the cleanest naturalistic-prompt
measurement of unfaithful CoT to date.%
\footnote{Excerpt status:
\texttt{kb/excerpts/arcuschin2025-cot-wild.md} is abstract-only as of
2026-05; \deepencite{arcuschin2025-cot-wild abstract,
citation-pending}} The methodological contribution is to drop the
artificial-bias scaffolding that earlier faithfulness work
(e.g., adding a known-misleading hint to the prompt) used to elicit
unfaithful traces, and instead measure rates on naturalistic, unbiased
prompts — the regime production users actually inhabit. The headline
per-model rationalisation rates: \textbf{GPT-4o-mini 13\%},
\textbf{Haiku 3.5 7\%} on production non-thinking models;
\textbf{Sonnet 3.7 thinking 0.04\%}, \textbf{Gemini 2.5 Pro 0.14\%},
\textbf{R1 0.37\%} on thinking-on configurations
\citep[abstract]{arcuschin2025-cot-wild}.%
\footnote{KB note:
\texttt{kb/notes/reasoning/chain-of-thought.md\#4.1}.} The
non-thinking-vs-thinking gap is roughly two orders of magnitude in
measured unfaithfulness rate and is the empirical anchor of the
faithfulness scaling question \cref{sec:faith-mehta} returns to.

A second contribution is the diagnosis of \emph{unfaithful illogical
shortcuts} on hard math problems: the model produces a long, plausibly
structured trace whose intermediate steps are individually defensible
but whose chain does not in fact derive the emitted answer; the answer
appears to be reached by a separate route the trace does not document.
\citep[\S 4]{arcuschin2025-cot-wild}.%
\footnote{\deepencite{arcuschin2025-cot-wild §4, citation-pending}}
This is the failure mode that breaks the simplest reading of
\cref{eq:cot-decode}: the conditional $y \sim p_\theta(\cdot \mid x,
z)$ is well-defined, but the de-facto computation of $y$ does not
flow through the trace tokens.

\subsection{Shen et al.\ FaithCoT-Bench: instance-level measurement}
\label{sec:faith-bench}

\citet{shen2025-faithcot-bench} promote faithfulness measurement from a
per-distribution rate to a \textbf{per-instance} discriminative
problem.%
\footnote{Excerpt status:
\texttt{kb/excerpts/shen2025-faithcot-bench.md} is abstract-only as of
2026-05; \deepencite{shen2025-faithcot-bench abstract,
citation-pending}} Their FINE-CoT dataset contributes
$> 1{,}000$ expert-annotated trajectories from four LLMs across four
domains, with $> 300$ labelled unfaithful, against which they
benchmark eleven detectors covering three paradigms:
\begin{itemize}
\item \textbf{Counterfactual detectors} apply intervention maps
$\mathcal{I}$ from \cref{sec:faith-def} and measure the answer-distribution
shift; an unfaithful trace is one where content-changing interventions
fail to change $y$ \citep[\S 3]{shen2025-faithcot-bench}.
\item \textbf{Logit-based detectors} read internal model signals
(token-level entropies, attention to trace tokens vs.\ prompt tokens)
to predict per-instance unfaithfulness without explicit intervention
\citep[\S 3]{shen2025-faithcot-bench}.
\item \textbf{LLM-as-judge detectors} present the trajectory to a
separate language model and ask whether the trace's stated reasoning
in fact derives the emitted answer
\citep[\S 3]{shen2025-faithcot-bench}.
\end{itemize}
The headline empirical finding: detection is \textbf{harder in
knowledge-intensive domains and on more advanced models}
\citep[abstract]{shen2025-faithcot-bench}. The intuition is uncomfortable
for the field's defaults: the very models whose trace-dependent answer
quality is best are the ones whose unfaithfulness is hardest to detect
mechanistically, because their traces are surface-coherent enough that
LLM-as-judge fails and their internal logit signal does not visibly
collapse on the unfaithful instances.

A second methodological point of FaithCoT-Bench is that
\emph{instance}-level measurement, not aggregate rate, is the
load-bearing object. A model with average rationalisation rate $0.14\%$
\citep[\S 4]{mehta2026-faithfulness-scaling} could still be unfaithful
on the specific subpopulation of prompts where the user actually depends
on trace transparency (high-stakes decisions, monitorable-channel
deployments — Paper~5's alignment surface). Aggregate measurement
underweights tail risk.

\subsection{Mehta et al.: does inference scaling improve faithfulness?}
\label{sec:faith-mehta}

\citet{mehta2026-faithfulness-scaling} ask the most consequential 2026
question for the paper's thesis: does \emph{scaling} the inference-time
compute axis (\cref{sec:test-time-compute}) — specifically,
self-consistency at $N > 1$ — \emph{also} improve faithfulness, or are
accuracy gains and faithfulness gains decoupled?%
\footnote{Excerpt status:
\texttt{kb/excerpts/mehta2026-faithfulness-scaling.md} is abstract-only
as of 2026-05; \deepencite{mehta2026-faithfulness-scaling abstract,
citation-pending}} The study runs four frontier models
(\textbf{GPT-5.2, Claude Opus 4.5, Gemini-3-flash-preview,
DeepSeek-v3.2}) on $100$ GSM8K problems at varying self-consistency
budgets and measures both task accuracy and a per-trajectory
faithfulness score in parallel.

The headline result is non-monotonic: \textbf{inference scaling does
not uniformly improve faithfulness}
\citep[abstract]{mehta2026-faithfulness-scaling}. The Claude Opus 4.5
configuration provides the cleanest illustration: accuracy
\emph{decreases} from $78\% \to 74.3\%$ while measured faithfulness
\emph{rises} from $0.270 \to 0.891$ at
$N = 5$.\footnote{KB index:
\texttt{kb/index/papers.json} — abstract.
\deepencite{mehta2026-faithfulness-scaling §4 (or wherever the per-model
table lives in the PDF; section anchor pending population), citation-pending}}
For \cref{eq:cot-decode}'s reading of self-consistency this is a
puzzle: more samples should average out per-sample noise in the
\emph{trace} channel, leaving the conditional $p_\theta(y \mid x, z)$
better-determined; but at the same time, the votes are over the
\emph{answer} channel, and the relationship between trace-channel
quality and answer-channel quality is exactly what faithfulness
measurement is testing. Mehta et al.'s finding is that the two channels
move differently per model.

The operational implication: \textbf{self-consistency must be tested
per-model before deployment}
\citep[abstract]{mehta2026-faithfulness-scaling} for any application
that depends on traces being a faithful read of the model's
computation. The framework's thesis — that inference compute, search,
and verification are co-designed levers (\cref{sec:test-time-compute,sec:process-reward-models})
— picks up a fourth operational variable here: the
\emph{interpretability cost} of the inference-compute leg, which the
2026 evidence reports is not a free externality of accuracy.

\subsection{The contradiction with reasoning-training framings}
\label{sec:faith-contradiction}

The clean read of \citet{deepseek-r1}'s GRPO outcome — long traces
emerge from RLVR, accuracy rises in lockstep with trace length,
truncating the trace destroys the gain — is that the trace is a
\emph{causal substrate}: the model is doing genuine compute on the
trace channel and the answer is being read off the result. The
Arcuschin and Mehta evidence complicates this read at two scales.

\begin{contradiction}
\textbf{Reasoning-training as causal-trace evidence vs.\ unfaithful-trace
evidence.} \citet[\S 2.2]{deepseek-r1} report that RLVR training grows
average response length from $\sim 200$ to $\sim 10{,}000$ tokens with
AIME 2024 pass\@1 climbing $15.6\% \to 77.9\%$,%
\footnote{KB excerpt:
\texttt{kb/excerpts/deepseek-r1.md\#sec-2-2-reward}.} a correlation
that the simplest reading attributes to the trace's content driving
the answer's correctness. Yet \citet{arcuschin2025-cot-wild} find R1
itself exhibits a measurable rationalisation rate of $0.37\%$
\citep[abstract]{arcuschin2025-cot-wild} on naturalistic prompts, and
the unfaithful-illogical-shortcut failure mode of
\cref{sec:faith-arcuschin} is a regime where trace tokens are emitted
in the right format but the answer is not in fact a function of them.
\citet{mehta2026-faithfulness-scaling}'s scaling-of-faithfulness result
sharpens the tension: one frontier model's accuracy drops while its
faithfulness rises under self-consistency, a pattern incompatible with
the simple reading that traces are a causal substrate at all scales.
The 2026 picture is regime-dependent and the question is unresolved;
\cref{sec:contradictions} catalogues this alongside the other live
faithfulness contradictions of \cref{sec:faith-arcuschin,sec:faith-mehta}.
\end{contradiction}

A second sub-tension within the empirical wave itself: the per-rate
gap of two orders of magnitude between non-thinking and thinking
configurations \citep[abstract]{arcuschin2025-cot-wild} cannot be
cleanly attributed to long-CoT SFT, RLVR post-training, or general
scale by published 2025--2026 ablations
\citep[\S 4]{mehta2026-faithfulness-scaling}; the controlled
ablation that varies one of these factors at fixed others has not been
run at frontier scale.%
\footnote{KB note:
\texttt{kb/notes/reasoning/chain-of-thought.md\#4.1}.} The reason
matters: if the gap is RLVR-driven, it is benchmarkable and tunable
through the training pipeline of \cref{sec:rlvr}; if it is scale-driven,
faithfulness is on a separate scaling curve from accuracy, and Mehta et
al.'s non-monotonic per-model results are the first signal of where the
two curves cross.

\subsection{Cross-paper consequences: probe surface and monitor channel}
\label{sec:faith-cross-paper}

The faithfulness question lands twice in later papers, and §11's
resolution shapes both.

\paragraph{Paper~4 — CoT as a probe surface.}
The interpretability programme of Paper~4 reads internal computation
through several methods: logit lenses, linear probes, sparse
autoencoders, activation patching, and circuit tracing. CoT traces are
the \emph{output-side} probe surface for that programme — they are
human-readable, token-aligned with the residual stream, and stable
under repeated decoding. The faithfulness question gates this use:
\textbf{if traces are unfaithful, the probe surface is decoupled from
the underlying computation}. Paper~4's §interpretability returns to
this section as a precondition for treating CoT-as-probe as a primary
methodology, and uses Mehta et al.'s per-model variation to argue that
probe-surface validity must be re-established model-by-model rather
than assumed from the family.

\paragraph{Paper~5 — CoT as an alignment monitor.}
The alignment programme of Paper~5 treats the CoT channel as a
\emph{monitorable surface}: detectors that read traces look for
warning signs of misalignment, scheming, sycophancy, and
alignment-faking before the answer is committed. The
canonical 2024 result on alignment-faking
\citep{greenblatt2024-alignment-faking} demonstrates that a model can
behave correctly on observable outputs while internally pursuing a
divergent objective; if the trace channel is unfaithful in
\cref{eq:faith-causal}'s sense, this monitoring approach has a blind
spot exactly where it most needs to see. Paper~5's §alignment cites
forward to this section as the empirical bound on CoT-as-monitor: a
$0.04$--$0.37\%$ unfaithfulness rate on thinking models
\citep[abstract]{arcuschin2025-cot-wild} is small in aggregate but
consequential at scale, and the per-instance harder-on-advanced-models
finding of \citet{shen2025-faithcot-bench} is the most direct warning
that the channel degrades exactly in the regime where alignment
monitoring matters most.

The shared structural point across both papers: \emph{aggregate
unfaithfulness rates do not bound deployment risk}. The instance-level
distribution matters, and the 2025--2026 evidence is that the unfaithful
instances cluster in knowledge-intensive domains and on advanced models
\citep[abstract]{shen2025-faithcot-bench} — the same regime where
downstream uses of the trace channel are most consequential.

\subsection{Forward link: does the verifier signal still matter?}
\label{sec:faith-forward}

The structural adjacency to \cref{sec:prm-survival} is direct. Both
sections ask the same shape of question: \emph{has the post-2024
training pipeline absorbed a function the field used to assign to an
external module?} For verifiers (PRMs), the question is whether
process-reward signal still earns its inference and training cost once
RLVR-trained policies internalise step-level verification. For traces,
the question is whether the CoT channel still earns its interpretability
and monitoring premium once the policy has been RL-trained to optimise
the answer-channel directly. The 2026 evidence, in both cases, is that
the absorption is partial and regime-dependent: PRMs continue to add
value at inference time (\cref{sec:prm-thinkprm,sec:prm-survival}) while
their training-time value shrinks, and traces continue to drive answers
on the on-distribution majority of instances while their reliability
degrades exactly on the harder, higher-stakes tail.
\Cref{sec:contradictions} catalogues both contradictions as twin live
frontiers, and the closing thesis of the paper —
\cref{sec:test-time-compute}'s triangle is co-designed in 2026 — must
inherit the qualifier that two of its three legs are operationally
sound but interpretively contested. \Cref{sec:prm-survival} carries the
verifier half of the question forward; this section's evidence is the
trace half.
